\documentclass{article}
\usepackage{graphicx}
\usepackage[margin=1in]{geometry}
\usepackage{booktabs}
\usepackage{longtable}
\usepackage{siunitx}
\usepackage{caption}
\usepackage{array}
\usepackage{hyperref}

\sisetup{
  round-mode=places,
  round-precision=3,
  group-separator={,},
  group-minimum-digits=4
}

\title{Applied Regression Homework 1 225040015}
\author{Dai Xunlian\\Student ID: 225040015\\225040015@link.cuhk.edu.cn}
\date{January 2026}

\begin{document}
\maketitle

\section{Introduction}
This report presents the conceptual answers and applied analyses for Homework 1.
All numerical results below are based on the provided computations.

\section{Conceptual}
\subsection*{1. (a)--(d)}
\begin{itemize}
  \item (a) \textbf{Flexible} is better: $n$ large, $p$ small, low overfitting risk, and flexible models can approximate the true relationship.
  \item (b) \textbf{Inflexible} is better: $p$ large, $n$ small, flexible models have high variance and overfit.
  \item (c) \textbf{Flexible} is better: the relationship is highly nonlinear and needs flexibility to capture curvature.
  \item (d) \textbf{Inflexible} is better: noise variance is high, so flexible models chase noise and increase variance.
\end{itemize}

\subsection*{3. Bias--Variance}
\paragraph{(a) Shape of curves (low $\to$ high flexibility).}
\begin{itemize}
  \item Bias$^2$: monotonically decreasing.
  \item Variance: monotonically increasing.
  \item Training error: monotonically decreasing.
  \item Test error: decreases first, then increases (U-shaped).
  \item Bayes error: flat (constant).
\end{itemize}

\paragraph{(b) Explanation.}
Increasing flexibility fits the training data more closely, lowering training error and bias.
However, it also makes the model more sensitive to sample fluctuations, increasing variance.
Test error combines bias$^2$, variance, and irreducible error, so it falls initially (bias reduction dominates)
and then rises when variance dominates. Bayes error depends only on noise and is unrelated to model choice.

\section{Applied 8: College Dataset}
\subsection*{(a) Read data}
\texttt{college <- read.csv("College.csv")}

\subsection*{(b) Set row names and drop school name column}
\texttt{rownames(college) <- college[,1]} \\
\texttt{college <- college[,-1]}

\subsection*{(c)(i) Summary statistics}
Table \ref{tab:college-summary} reports the minimum, quartiles, median, mean, and maximum of each numeric variable.

\setlength{\tabcolsep}{4pt}
\begin{longtable}{l S S S S S S}
\caption{Summary of numeric variables in the College dataset}\label{tab:college-summary}\\
\toprule
Variable & {Min} & {Q1} & {Median} & {Mean} & {Q3} & {Max} \\
\midrule
\endfirsthead
\toprule
Variable & {Min} & {Q1} & {Median} & {Mean} & {Q3} & {Max} \\
\midrule
\endhead
Apps & 81 & 776 & 1558 & 3001.64 & 3624 & 48094 \\
Accept & 72 & 604 & 1110 & 2018.8 & 2424 & 26330 \\
Enroll & 35 & 242 & 434 & 779.973 & 902 & 6392 \\
Top10perc & 1 & 15 & 23 & 27.5586 & 35 & 96 \\
Top25perc & 9 & 41 & 54 & 55.7967 & 69 & 100 \\
F.Undergrad & 139 & 992 & 1707 & 3699.91 & 4005 & 31643 \\
P.Undergrad & 1 & 95 & 353 & 855.299 & 967 & 21836 \\
Outstate & 2340 & 7320 & 9990 & 10440.7 & 12925 & 21700 \\
Room.Board & 1780 & 3597 & 4200 & 4357.53 & 5050 & 8124 \\
Books & 96 & 470 & 500 & 549.381 & 600 & 2340 \\
Personal & 250 & 850 & 1200 & 1340.64 & 1700 & 6800 \\
PhD & 8 & 62 & 75 & 72.6602 & 85 & 103 \\
Terminal & 24 & 71 & 82 & 79.7027 & 92 & 100 \\
S.F.Ratio & 2.5 & 11.5 & 13.6 & 14.0897 & 16.5 & 39.8 \\
perc.alumni & 0 & 13 & 21 & 22.7439 & 31 & 64 \\
Expend & 3186 & 6751 & 8377 & 9660.17 & 10830 & 56233 \\
Grad.Rate & 10 & 53 & 65 & 65.4633 & 78 & 118 \\
\bottomrule
\end{longtable}

\noindent Private counts: Yes = 565, No = 212.

\subsection*{(c)(ii) Pair plots (first 10 variables)}
From the scatterplot matrix:
\begin{itemize}
  \item \texttt{Apps--Accept--Enroll--F.Undergrad} show strong positive correlations.
  \item \texttt{Top10perc--Top25perc} shows a strong positive correlation.
  \item \texttt{Outstate} is positively related to \texttt{Expend} and \texttt{Top10perc}.
\end{itemize}

\subsection*{(c)(iii) Outstate vs Private}
Private schools have clearly higher out-of-state tuition.
Table \ref{tab:college-outstate} summarizes \texttt{Outstate} by group.

\begin{table}[h]
\centering
\caption{Outstate tuition by group}
\label{tab:college-outstate}
\begin{tabular}{l S S S S S}
\toprule
Group & {Min} & {Q1} & {Median} & {Q3} & {Max} \\
\midrule
Private=Yes & 2340 & 9100 & 11200 & 13970 & 21700 \\
Private=No & 2580 & 5366 & 6609 & 7844 & 15732 \\
Elite=Yes & 5224 & 12219 & 16950 & 18411.5 & 20100 \\
Elite=No & 2340 & 7070 & 9556 & 12155 & 21700 \\
\bottomrule
\end{tabular}
\end{table}

\subsection*{(c)(iv) Elite indicator}
Define \texttt{Elite} by \texttt{Top10perc} $>$ 50.
Counts: Yes = 78, No = 699.
Elite schools have much higher \texttt{Outstate} tuition (median 16{,}950 vs 9{,}556).

\subsection*{(c)(v) Histograms}
The distributions of \texttt{Apps}, \texttt{Enroll}, \texttt{F.Undergrad}, and \texttt{Expend} are strongly right-skewed.
\texttt{Outstate} is moderately right-skewed, and \texttt{Grad.Rate} has a few values above 100.

\subsection*{(c)(vi) Summary}
\begin{itemize}
  \item Private and Elite schools tend to have higher out-of-state tuition.
  \item \texttt{Apps--Accept--Enroll--F.Undergrad} form the strongest size-related cluster.
  \item \texttt{Top10perc} and \texttt{Top25perc} are highly correlated, indicating consistent student quality measures.
  \item \texttt{Expend} is positively related to \texttt{Outstate} and \texttt{Top10perc}, suggesting resources and student quality move together.
\end{itemize}

\subsection*{Correlation highlights}
Table \ref{tab:college-corr} lists several variable pairs with high absolute correlations.

\begin{table}[h]
\centering
\caption{Highly correlated variable pairs in the College dataset}
\label{tab:college-corr}
\begin{tabular}{l l S}
\toprule
Variable 1 & Variable 2 & {Correlation} \\
\midrule
Enroll & F.Undergrad & 0.9646 \\
Apps & Accept & 0.9435 \\
Accept & Enroll & 0.9116 \\
Top10perc & Top25perc & 0.8920 \\
PhD & Terminal & 0.8496 \\
Outstate & Expend & 0.6728 \\
S.F.Ratio & Expend & -0.5838 \\
\bottomrule
\end{tabular}
\end{table}

Enrollment scale variables are strongly positively correlated. Spending (\texttt{Expend}) is positively correlated
with \texttt{Outstate}, while student-faculty ratio is negatively correlated with spending.

\section{Applied 10: Boston Dataset}
\subsection*{(a) Basic information}
\texttt{dim(Boston)} = 506 rows and 14 columns (13 predictors + \texttt{medv}).

\subsection*{(b) Pairwise scatterplots}
Key relationships from the scatterplot matrix:
\begin{itemize}
  \item \texttt{rad--tax} is extremely positively correlated (about 0.91).
  \item \texttt{indus--nox} is positively correlated; \texttt{dis--nox} is negatively correlated.
  \item \texttt{rm--medv} is positively correlated; \texttt{lstat--medv} is strongly negatively correlated.
  \item \texttt{age--dis} is negatively correlated.
\end{itemize}

\subsection*{(c) Correlation with crime rate}
Variables most correlated with \texttt{crim} (by absolute value):

\begin{table}[h]
\centering
\caption{Correlation with \texttt{crim} in the Boston dataset}
\label{tab:boston-corr}
\begin{tabular}{l S}
\toprule
Variable & {Correlation} \\
\midrule
rad & 0.626 \\
tax & 0.583 \\
lstat & 0.456 \\
nox & 0.421 \\
indus & 0.407 \\
medv & -0.388 \\
b & -0.385 \\
dis & -0.380 \\
\bottomrule
\end{tabular}
\end{table}

\subsection*{(d) Extremes and ranges}
Maximum \texttt{crim} = 88.9762 (row 381), maximum \texttt{tax} = 711 (row 489),
and maximum \texttt{ptratio} = 22 (row 355).
Table \ref{tab:boston-range} summarizes each variable's range.

\begin{table}[h]
\centering
\caption{Variable ranges in the Boston dataset}
\label{tab:boston-range}
\begin{tabular}{l S S}
\toprule
Variable & {Min} & {Max} \\
\midrule
crim & 0.00632 & 88.9762 \\
zn & 0 & 100 \\
indus & 0.46 & 27.74 \\
chas & 0 & 1 \\
nox & 0.385 & 0.871 \\
rm & 3.561 & 8.78 \\
age & 2.9 & 100 \\
dis & 1.1296 & 12.1265 \\
rad & 1 & 24 \\
tax & 187 & 711 \\
ptratio & 12.6 & 22 \\
b & 0.32 & 396.9 \\
lstat & 1.73 & 37.97 \\
medv & 5 & 50 \\
\bottomrule
\end{tabular}
\end{table}

\subsection*{(e) Charles River}
\texttt{sum(Boston\$chas == 1)} = 35.

\subsection*{(f) Median pupil--teacher ratio}
\texttt{median(Boston\$ptratio)} = 19.05.

\subsection*{(g) Minimum-\texttt{medv} tract}
The minimum \texttt{medv} is 5.0 at row 399. Its characteristics are listed in Table \ref{tab:boston-minmedv}.

\begin{table}[h]
\centering
\caption{Variables for the minimum-\texttt{medv} tract}
\label{tab:boston-minmedv}
\begin{tabular}{l S}
\toprule
Variable & {Value} \\
\midrule
crim & 38.3518 \\
zn & 0 \\
indus & 18.1 \\
chas & 0 \\
nox & 0.693 \\
rm & 5.453 \\
age & 100 \\
dis & 1.4896 \\
rad & 24 \\
tax & 666 \\
ptratio & 20.2 \\
b & 396.9 \\
lstat & 30.59 \\
medv & 5 \\
\bottomrule
\end{tabular}
\end{table}

\noindent This tract has high \texttt{crim}, \texttt{nox}, \texttt{tax}, \texttt{ptratio}, and \texttt{lstat},
with low \texttt{rm} and \texttt{dis}, consistent with older, polluted, and less affluent areas.

\subsection*{(h) Counts of large houses}
\texttt{sum(Boston\$rm > 7)} = 64 and \texttt{sum(Boston\$rm > 8)} = 13.
Tracts with \texttt{rm} $>$ 8 tend to have higher housing prices and lower \texttt{lstat}.

\section{Conclusion}
\begin{itemize}
  \item Conceptual results highlight the bias--variance tradeoff and when flexible vs.\ inflexible models are preferred.
  \item In the College data, private and Elite schools show higher out-of-state tuition, and enrollment-related
  variables are strongly correlated.
  \item In the Boston data, crime rate is positively related to accessibility and tax rate, and negatively
  related to housing prices and distance; the minimum-\texttt{medv} tract exhibits extreme values.
\end{itemize}

\end{document}
